% Generated by scripts/tikz_reviews.py from the figure manifests.
\ReviewFigure{Entropy on the probability simplex}
{figures/1-shannon/entropy-extrema.pdf}
{figures/tikz/coding-entropy-extrema.pdf}
{The entropy lemma states that point masses minimize entropy and the uniform distribution maximizes it.}
{Interpretation to check: The left sketch has unclear level-curve labels; it is recast as the exact binary entropy graph on p=(t,1-t). The right panel retains the three-symbol simplex, with numerically computed entropy contours at 1.05, 1.30, and 1.50 bits. Cropped prose in the scan is omitted.}
{coding-entropy-extrema}
\ReviewFigure{Kraft inequality: disjoint descendant blocks}
{figures/1-shannon/kraft-ineq-1.png}
{figures/tikz/coding-kraft-necessity.pdf}
{The necessity proof completes a prefix tree to depth m and counts the descendants of each codeword.}
{The code is exactly the one specified in the caption: (0,10,110,111). The corresponding descendant blocks contain 4, 2, 1, and 1 leaves of the full depth-three tree. Highlighted circles are codeword roots; the colored leaf blocks are disjoint.}
{coding-kraft-necessity}
\ReviewFigure{Kraft inequality: packing prescribed lengths}
{figures/1-shannon/kraft-ineq-2.png}
{figures/tikz/coding-kraft-sufficiency.pdf}
{The converse proof packs dyadically aligned blocks in nonincreasing order of size.}
{For lengths (1,2,3), the blocks have 4, 2, and 1 leaves. Packing them from the left gives codewords (0,10,110), leaving 111 unused. The tree demonstrates both alignment and the inequality 7/8 <= 1; it does not assume equality.}
{coding-kraft-sufficiency}
